
\section{Monetary Theory}

\begin{frame}{Motivation}
    \mycomment{
    The world is the only closed economy we have. So there is a field to study Global Macroeconomics. 
    
    An intermediate step is to look at the Open Economy, which applies established macroeconomic models to the economy and adds the Rest of the World as given and outside the realm of politics. 
    }
    \mytext{
    \begin{itemize}
        \item So far, only goods and services were traded
        \item so far: \textit{microeconomics}
        \item How about the passive side of the balance sheet?
        \item How about trading financial assets, capital, and debt?
        \item What about monetary effects, exchange rates, inflation?
        \item What about economic growth and development?
        % \item $\Rightarrow$ Macroeconomics Global Macroeconomics
    \end{itemize}
    }
\end{frame}


\begin{frame}{Recap: Balance of Payments}
    \showmytikz{../shared/tikz/tikz-boptikz}
    \mycomment{
    % In 2025 the EU exported \euro{2,584} Trillion (!) in goods and imported \euro{2,434}. The difference - 150 Trillion - 
    % This accounts for 30.7\% of world exports and 28.5\% of world imports.
    
    % The EU exported \euro{1.427} Billion (!) in services, and exported services for \euro{1.274} Billion. 

    For the EU BoP see \cite{eubop}

    For the German BoP see \cite{bundesbankbop2024}
    }
\end{frame}

\begin{frame}{Recap: Balance of Payments}
    \showmytikz{../shared/tikz/tikz-bop-entries}
\end{frame}

% \begin{frame}{Balance of Payments (detailed)}
%     \showmytikz{../shared/tikz/tikz-bop}
% \end{frame}

\begin{frame}{Debt / Liability / Money}
    \mytext{
    \begin{itemize}
        \item If I give my daughter only 100 Euro a month...
        \item ... but she buys stuff for 200 Euro a month...
        \item there must be some other source of income!
        \item ... so either she sells something else...
        \item ... or she lives of borrowed money.
        \item Even if she steals, she will accumulate liabilities!
    \end{itemize}
    }
\end{frame}

\begin{frame}{Why Borrow or Lend Money?}
    \mytext{
    \begin{itemize}
        \item Maybe she lives of capital transfer? Like her grandmother sponsoring her?
        \item Maybe she is over-spending off her credit card?
        \item Maybe someone lends her money?
        \item So, she issues bonds?
        \item What are the costs of that borrowed money?
    \end{itemize}
    }
\end{frame}

\begin{frame}{The Costs of Money: Interest Rates}
    \mytext{
    \begin{itemize}
        \item The cost of money is called interest.
        \item Interest is typically expressed as interest~rate~($i$).
        \item There is also a (hidden) cost of money: inflation~$(\pi)$.
        \item Internationally, there is also (the chance of) devaluation (or revaluation) of currency~($\Delta Exp$).
    \end{itemize}
    }
\end{frame}

% \begin{frame}{Money}
%     \mycomment{
%     Money is not a commodity like any other. But what exactly is money, where does it come from, what constitutes its value and what guarantees that it keeps this value?
%     }
%     \mytext{
%     \begin{itemize}
%         \item There can be debt without money
%         \item However, most debt will be monetary
%         \item We must think about money
%     \end{itemize}
%     }
% \end{frame}

\begin{frame}{History of Money}
    \mycomment{
    For a long time, precious metals served as money. Little was understood that the actual function of money is different to that of any other good, that it is actually not a "good" in the same way that all other commodities are.

    The early story of money is a story of divergence of the actual face value of coins relative to the metal price, leading to the eventual decoupling of the two.
    }
\mytext{
    \begin{itemize}
        \item Money - a commodity?
        \begin{itemize}
            \item gold, gold standard 
            \item silver ...
            \item Gresham's law
            \item paper money?
            \item digital money?
        \end{itemize}
    \end{itemize}
}
\end{frame}

\begin{frame}{History of Money}
    \mycomment{
    Isaak \noun{Newton} became Head of the Royal Mint, and he wanted to \textit{discourage} the use of gold as currency in favour of silver. ("Sterling" is actually a silver term.) He thus decreed that while before, the gold:silver ratio was about 1:20, it would now be exchanged 1:21, believing this would disfavour gold.

    But the effect was that foreigners brought in gold, that was now more expensive than silver, while silver left the country. The "gold standard" was acknowledged almost a century later.
    }
    \mytext{
    \begin{itemize}
        \item Bank of England issues paper money: 1694
        \item Pound Sterling gold standard: 1717, 1821
        \item Banque de France issues paper money: 1800
        \item Greenback: American Civil War United States Notes
        \item Federal Reserve System FED 1913
        \item Reichsmark gold standard until 1914
        \item German hyperinflation
        \item Russian attempts to abolish money 1917-1921
        \item Bretton Woods system until 1973
        \item Euro 2002
    \end{itemize}
    }
\end{frame}



\begin{frame}{History of Money}
    \mycomment{
    In addition to central banks (or mints) issuing currency, the banking sector developed. This allows the creation of "book money" and to conduct payments without the actual shipping of currency. Using the banking system, merchants could now ship wares back and forth and square their accounts without the perilous transportation of precious metals in the opposite direction.

    By doing so, banks can indeed create book money. This must be limited by some required guarantees and underlying reserves to preserve bank runs. The supervision of banks is one of the chief responsibilities of the central banks.
    }
    \mytext{
    \begin{itemize}
        \item Sconto, conto, credit, debit: Italian
        \item Giro system 17th century, Marc Banco (Hamburg 1619, Nurnberg, Venice...)
        \item Banking system / double book keeping: Florence
        \item Securities, bonds, stocks (Wertpapier)
        \item "check": English
    \end{itemize}
    }
\end{frame}


\begin{frame}{History of Money}
    \mycomment{
    Along with the development of new means of payments came the development of other financial institutions, where debt in the form of shares and stocks can be traded. This came along with new forms of companies, financed by a multitude of shareholders. Those companies are also liable only up to the underlaying investment (limited company), and are their own legal personality, and without any person holding liability.
    }
    \mytext{
    \begin{itemize}
        \item First stock exchange Brügge 1276
        \item Paris 1572
        \item Hamburg 1558, Berlin 1685, Bremen 1682
        \item London 1773, New York 1792
    \end{itemize}
    }
\end{frame}


\begin{frame}{History of Money}
    \mycomment{
    The history of money is not over, as we keep developing new means of payment. Digitalisation and the invention of strong encryption makes new forms of payment possible. The underlying principle though, that states and their request for taxes (legal tender) are the ultimate guarantors of the value of money, remains unchanged.
    }
    \mytext{
    \begin{itemize}
        \item Credit cards: Diners Club 1950, 1977 VISA
        \item ATMs since 1960, Online Banking since 2000s
        \item Debit cards / Maestro, VISA debit
        \item SWIFT 1973
        \item Crypto: DES, Triple DES, AES since 1977
        \item Bitcoin 2008, Ethereum
        \item CBDC (Central bank digital currency)?
        \item Quantum computer?
    \end{itemize}
    }
\end{frame}

\begin{frame}{FOREX global currency exchange}
\mycomment{
    FOREX measures only the directly exchanged currency. 
    
    The USD plays a much more important role as unit of account in the vastly larger currency swap markets \cite{blustein2025}.

    In addition, most commodities are traded in USD. The role of the USD is thus not just (not so much) that of a reserve currency, but that of (international) unit of account.
}
\mytable{
    \begin{tabular}{lr}
        \toprule
        Currency & Mio USD / day \\
        \midrule
         USD &  6,638,790 \\
         EUR & 2,292,444 \\
         JPY  & 1,253,236 \\
         GBP & 968,139 \\
         CNY & 526,231 \\
         CAD & 465,776 \\
         CHF & 389,700 \\
         other & 1,999,093 \\
         \bottomrule
    \end{tabular}
    
    According to BIS, 2022 
}

\end{frame}

\begin{frame}{Currency Reserves of Central Banks, 2023 (IMF)}
\mycomment{
    Currency reserves are sometimes held in actual cash (bank notes), but mostly in government bonds. 
    
    The most important of these are USTs (United States Treasuries), JSTs (Japanese State Treasuries), Gilts (UK treasuries), and various European bonds, often summarised as EU10Y. The most common form is a 10 year bond.

    In addition, central banks hold other forms of reserve, the most importantly gold, which is traded (denominated) in USD. Put together, central banks hold a little more than half their reserves in USD or USD dominated assets (such as USTs).
}
\mytable{

    \begin{tabular}{lr}
    \toprule
    Currency & Bio USD \\
    \midrule
         USD & 6,687  \\
         EUR & 2,287 \\ 
         JPY & 652 \\
         GBP & 553 \\
         CNY & 261 \\
        CAD & 295 \\
        SFR & 26 \\
        other & 883 \\
    \bottomrule
             
    \end{tabular}
}
\end{frame}


\begin{frame}{Global Financial System}
\mycomment{
    The global financial system is a "non system" since the breakdown of fixed exchange rates in the 1970ies. Today, a number of currencies compete. This also means that there is a competition for credibility and stability between currencies. This exercises competitive pressure on interest rates, tax rates, and inflation rates. The position of the US Dollar is challenged, but it remains the most important currency in many aspects, and that is unlikely to quickly change.
}
\mytext{
    \begin{itemize}
        \item Central banks, independent (?)
        \item M2 cash + short term deposits
        \item Flexible exchange rates (really?)
        \item free movement of capital since 1990ies
        \item Bank for International Settlements BIS
        \item Basel Committee on Banking Supervision
        \item IMF Lender of Last Resort (?)
    \end{itemize}
}
\end{frame}

\begin{frame}{The Four Functions of Money}
\mycomment{
    For a currency to be accepted money, it must serve four functions. Other commodities like gold and silver may be highly liquid, as they can be exchanged \textit{for money} easily. But they are not money still as long as there is \textit{real money} that serves these functions better.
    
    Right after the Second World War, my father managed to exchange a silver spoon for a quart of milk. But he could have bought that quart easier (and probably a larger quantity) for US dollars, had he had some.
}
\mytext{
    The four functions of money are:
    \begin{enumerate}
        \item medium of exchange
        \item unit of account
        \item storage of value
        \item standard of deferred payment
    \end{enumerate}
}
\end{frame}

\begin{frame}{Currency}
\mycomment{
It is important to remember the general impossibility to compare the use value (that is individual) to the exchange value (that is determined on markets) of any good. Money has no use value at all but is all exchange value, which is its strength.

Notice that the exchange value is at parts dependent on the costs to produce a good---that is, in competitive markets. In non-competitive markets, the value can be higher than the costs to produce.
}
\mytext{
    There is a fundamental difference between
    \begin{itemize}
        \item use value, and
        \item exchange value.
        % \item intrinsic value (strike price vs. underlying)
    \end{itemize}

    Money has no use value.
}
\end{frame}

\begin{frame}{Currency}
\mycomment{
States guarantee the value of currency, and they actually force their subjects to use it. States hold a monopoly on money. Nobody except the central bank can issue currency. And banks creating book money are under supervision of the central banks.
}
\mytext{
    \begin{itemize}
        \item Obligation to accept (for debt settlement)
        \item Paying taxes etc. in that currency
        \item fiat money: no intrinsic value
        \item legal tender
    \end{itemize}
}
\end{frame}

\begin{frame}{Currency}
\mytext{
    For a currency to work, it must be accepted, that means:
    \begin{itemize}
        \item stability
        \item scarcity
        \item convertibility
    \end{itemize}
}
\mycomment{
    Currencies compete with one another in an open economy. However, some countries---most notable China---have opted out of full convertibility. This means that the "currency of the people" (Renminbi, with its unit "Yuan" meaning "unit") is stable and scarce, but not convertible. Other currencies are convertible, but often not stable. There is much to be said about optimum currency areas and whether or not a currency should be pegged to another, to a bundle, or free-float.
}
\end{frame}


\begin{frame}{Quantitative Theory of Money}
\showmytikz[1,2]{../shared/tikz/tikz-quantitytheoryofmoney}
\mycomment{
    Imagine a two-sector economy, with only fishermen and coconut farmers.

    The (exchange) value of fish in coconut will depend entirely on how many fish and how many coconut are brought to the market (per market day). This in turn will depend on preferences of the consumers, and on the effort needed for fishing and harvesting. 

    Imagine a seasonal shortage of coconut relative to plenty of fish. Then some fishermen will change profession and start harvesting coconut. 
    
    If labour is mobile (and all that counts), then costs to fish and harvest should equal out.

    This is a world of exchange without money, which may work if there are only fish and coconut. 
    
    As commerce expands beyond fish and coconut, a common denominator will become desirable. 
}
\end{frame}

\begin{frame}{Quantitative Theory of Money}
\showmytikz[1,2,3,4]{../shared/tikz/tikz-quantitytheoryofmoney}
\mycomment{
        Now our islanders agree using some sea shells (of an extinct species) as medium of exchange. 
        
        Farmers buy fish with shells, and fishermen with shells buy coconut. 
    
        The value of the shells (money) depends only on two factors:
        \begin{itemize}
            \item the speed of circulation
            \item the quantity of shells in circulation
        \end{itemize}
}
\end{frame}

\begin{frame}{Quantitative Theory of Money}
\showmytikz[1,3]{../shared/tikz/tikz-quantitytheoryofmoney}
\mycomment{
        For the system (or better: this model!) to work, the shells must retain circulation. For they do not have any use value, fishermen would not accept shells without farmers accepting shells.
        
        So why shells? And where do these shells come from? One quality the shells have is that nobody can create them at will. Really? But we can create debt, can we not?  
}
\end{frame}

\begin{frame}{Quantitative Theory of Money}
\showmytikz[1,3,5]{../shared/tikz/tikz-quantitytheoryofmoney}
\mycomment{
        In our model world, neither $A$ nor $B$ ever invest in new fishing nets or ladders. If a farmer wants to take time off to make a new ladder instead of harvesting, they would need to get some credit to not starve while they do this. The fishermen would have to supply shells as a loan (and then fish in exchange for shells). Where would that credit come from, if the farmers need all their money for transactions?

        In other words: the quantitative theory of money, while absolutely worth knowing and considering, has its limits. It assumes there is a world without credit, without debt, without interest rate, without past or future, without inequality, without power and ownership, without permissions, competing currencies, etc.
}
\end{frame}


\begin{frame}{Theory of Money  I - Debit First}
    \mycomment{
    One theory---a popular, but disproven one---assumes that banks came into existence after money. That is, there was money in circulation (gold, silver etc.). Banks were founded to safeguard  gold and silver.

    We call this "debit first" theory. In this world, debit is everything; credit comes later, or remains an exception.
    }
    \mytable{
            \begin{tabular}{p{20ex}|p{20ex}}
                \multicolumn{2}{c}{My Business}\\
                 debit (assets) & credit (liabilities) \\\hline            
                 500 stuff & \\
                 \tikzmark{gold1} 500 gold  &  1000 equity\\\hline
            \end{tabular}
    
           \bigskip
    
            \begin{tabular}{p{20ex}|p{20ex}}
                \multicolumn{2}{c}{My Bank}\\
                 debit (assets) & credit (liabilities) \\\hline            
                 \tikzmark{gold2} 0 gold   & 0 bank notes \\\hline
            \end{tabular}
            \begin{tikzpicture}[overlay, remember picture]
                \draw[red,->,bend left=-90] (pic cs:gold1) to (pic cs:gold2);
            \end{tikzpicture}
    }
\end{frame}

\begin{frame}{Theory of Money  I - Debit First}
    \mycomment{
    According to this, banks would collect the "hard cash" (gold, silver etc.) and issue bank notes instead. 
    }
\mytable{
        \begin{tabular}{p{20ex}|p{20ex}}
            \multicolumn{2}{c}{My Business}\\
             debit (assets) & credit (liabilities) \\\hline            
             500 stuff &  \\ 
             500 bank notes \tikzmark{note2}  &  1000 equity\\\hline
        \end{tabular}

        \bigskip

        \begin{tabular}{p{20ex}|p{20ex}}
            \multicolumn{2}{c}{My Bank}\\
             debit (assets) & credit (liabilities) \\\hline            
             500 gold & \tikzmark{note1} 500 bank notes \\\hline
        \end{tabular}
        \begin{tikzpicture}[overlay, remember picture]
                \draw[red,->,out=180,in=0] (pic cs:note1) to (pic cs:note2);
        \end{tikzpicture}
}
\end{frame}

\begin{frame}{Theory of Money  I - Debit First}

\mycomment{
    In the next step, we can substitute the bank notes even further---by current accounts or digital money. Whether there is  an actual paper bill or an electronic proof, the idea is the same: The Theory of Money here is debit. 

    While this is a nice theory, it does not explain where money came from in the first place. We now know that the value of money has other origins.
}

\mytable{
        \begin{tabular}{p{20ex}|p{20ex}}
            \multicolumn{2}{c}{My Business}\\
             debit (assets) & credit (liabilities) \\\hline            
             500 stuff &  \\ 
             500 debit to bank &  1000 equity\\\hline
        \end{tabular}

        \bigskip

        \begin{tabular}{p{20ex}|p{20ex}}
            \multicolumn{2}{c}{My Bank}\\
             debit (assets) & credit (liabilities) \\\hline            
             500 gold & 500 credit from me \\\hline
        \end{tabular}
}        
\end{frame}



\begin{frame}{Theory of Money  II - Credit First}
    \mycomment{
    Another theory turns the argument upside down. Rather than "us" just having money out of nowhere, we begin with nothing. It is the banks that start supplying the economy with money. The bank creates money "ex nihilo". 

    This is a credit economy, so we call it "credit first", and it is actually more realistic.

    In the first step, the individual(s) have nothing; the bank prints some bank notes, and holds (assumingly) "something" to safeguard its value, and establish trust. 
    }
    \mytable{
            \begin{tabular}{p{20ex}|p{20ex}}
                \multicolumn{2}{c}{My Business}\\
                 debit (assets) & credit (liabilities) \\\hline            
                 0 stuff & \\
                 0 gold &  0 equity\\\hline
            \end{tabular}
    
            \bigskip
    
            \begin{tabular}{p{20ex}|p{20ex}}
                \multicolumn{2}{c}{My Bank}\\
                 debit (assets) & credit (liabilities) \\\hline            
                 500 gold (?) & 500 equity (?) \\\hline
            \end{tabular}
    }        
\end{frame}

\begin{frame}{Theory of Money  II - Credit First}
    \mycomment{
        In the second step, it lends these bank notes to the economy. 
    }
    \mytable{
            \begin{tabular}{p{20ex}|p{20ex}}
                \multicolumn{2}{c}{My Business}\\
                 debit (assets) & credit (liabilities) \\\hline            
                 0 stuff &  \tikzmark{cred2} 500 credit from bank \\
                 500 bank notes \tikzmark{notes2} &  0 equity\\\hline
            \end{tabular}
    
            \bigskip
    
            \begin{tabular}{p{20ex}|p{20ex}}
                \multicolumn{2}{c}{My Bank}\\
                 debit (assets) & credit (liabilities) \\\hline         
                 500 debit to me \tikzmark{cred1} & 500 equity (?) \\
                 500 gold (?) & \tikzmark{notes1}(500 bank notes) \\\hline
            \end{tabular}

        \begin{tikzpicture}[overlay, remember picture]
                \draw[red,->,out=180,in=0] (pic cs:notes1) to (pic cs:notes2);
                \draw[green,ultra thick,<->,>=latex,out=-180,in=0] (pic cs:cred2) to (pic cs:cred1);
        \end{tikzpicture}
    }        
\end{frame}


\begin{frame}{Theory of Money  II - Credit First}
\mycomment{
    Now the economic subjects can use this to start trading, invest into machinery, and produce. 
    
    They will at the end invest the money in a way that what is produced at the end of the period is larger than what there was before. Thus, money supply must grow in parallel to the growth of the economy. Also, to be able to grow, the economy needs cash supply.

    \begin{itemize}
        \item Gold is not needed
        \item creation of money "ex nihilo"
        \item but: \textbf{interest rates}
        \item normally, \textbf{interest rate = growth rate}
    \end{itemize}
}
\mytable{
        \begin{tabular}{p{20ex}|p{20ex}}
            \multicolumn{2}{c}{My Business}\\
             debit (assets) & credit (liabilities) \\\hline            
             \tikzmark{stuff} 500 investment\tikzmark{stuff2} & \tikzmark{creditfrombank}500 credit from bank\tikzmark{creditfrombank2} \\ 
             \tikzmark{noty} 0 bank notes &  0 equity \\\hline
        \end{tabular}

        \bigskip

        \begin{tabular}{p{20ex}|p{20ex}}
            \multicolumn{2}{c}{My Bank}\\
             debit (assets) & credit (liabilities) \\\hline         
             \tikzmark{debittome}500 debit to me &  1000 equity\\
             500 gold (?) &  0 banknotes  \\\hline
        \end{tabular}

        \begin{tikzpicture}[overlay, remember picture]
                \draw[red,->,out=180,in=180,looseness=2] (pic cs:noty) to (pic cs:stuff);
                \draw[red,->,out=0,in=180] (pic cs:creditfrombank2) to node[above] {\scriptsize debt service} (pic cs:debittome);
                \draw[->,red] (pic cs:stuff2) -- node[above]{\scriptsize interest} (pic cs:creditfrombank);
        \end{tikzpicture}

}        
\end{frame}

\begin{frame}{Banking}
    \showmytikz{../shared/tikz/tikz-bankingsector}
    \mycomment{
    The central bank lends money to banks on the condition of interest payments. The banks lend money to businesses, giving them credit, on the condition of interest payments.

    The businesses employ consumers as workers and labourers, produce goods more worth than the credit they have received, paying the salaries and the interest rates.

    The systems works so that the growth rate needs to relate to the interest rate. The entire system hinges on the central bank, setting the central interest rate.
    }
\end{frame}

\begin{frame}{Control Money Supply?}
    \mycomment{
        Since banks can create money out of nothing, there is a temptation to do so. But that undermines the stability of the issued currency through inflation, and undermines trust.

        Thus, banks are strictly regulated by the central bank, which is the only institution allowed to issue currency. Private banks can participate in the creation of money and value, but under strict oversight. 

        The independence of the central bank means that it serves the public interest best while it serves its own interests best.
    }
    \mytext{
        \begin{itemize}
            \item How to stop banks from printing money?
            \item How to control value of money?
            \item How to control inflation?
        \end{itemize}
    }
\end{frame}



\begin{frame}{M2 Money Supply}
\mycomment{
Central banks have a number of tools to control the amount of money in circulation. But what exactly counts as money?
}
\mytext{
    \begin{itemize}
        \item cash
        \item checking accounts
        \item savings accounts
        \item other (certificats of deposit)
        \item not retirement accounts and time deposits
    \end{itemize}

    Central banks typically target M2

    $\rightarrow$ to control inflation

    \bigskip

    Central banks main policy lever is the interest rate. In addition are the rules restricting money creating by commercial banks, such as reserve requirements.
    }
    
\end{frame}

\begin{frame}{Money Demand}
    \mytext{
    Recall the four functions:
        \begin{enumerate}
        \item medium of exchange
        \item unit of account
        \item storage of value
        \item standard of deferred payment
    \end{enumerate}

    \medskip
    
    \textbf{1-2 deal with current output} ($Y$)

    \textbf{3-4 deal with the future }
    }
    \mycomment{
    The nature of money demand is tricky. People will need more money as they exchange more, so the amount of money must grow with output to keep the value of money stable.

    The inter-temporal aspects of money must guarantee the trust into the currency in the long run. Thus the (real) interest rate should correspond to the growth rate of the output $Y$, often abbreviated $g$.
    }
\end{frame}


\begin{frame}{Nominal vs. Real Interest Rate}
    \mytext{

    Consider:

    \begin{itemize}
        \item Growth rate of GDP $g=\frac{dY}{dt}$
        \item Inflation $\pi=\frac{dP}{dt}$
        \item Real interest rate $r=i/\pi$ 
        \item Nominal interest rate $i$ 
    \end{itemize}    
    }
    \mycomment{
    It is not exactly clear what causes the economy to grow. 
    
    The most obvious is population growth: if output per person is stable, more persons will output more. 

    On top, technological progress, increased efficiency, and new energy sources lead to increased output per head, and thus to growth.

    Whether the interest rate reacts to growth or does stipulate growth is a matter of debate.
    
    }
\end{frame}


% \begin{frame}{Nature of Interest Rates}
% \mycomment{
% While it is clear that the price of money is the interest rate, it is less clear where this interest rate actually "comes from". Is it set voluntarily by the monopolist (that is, the issuer of currency, the central bank), or is it a market outcome?

% The answer is probably in the middle. Yes, it is set by the central bank, which acts like a monopolist. But it is also a market result, just like any monopolist price. The central bank reacts to a market, the money market. But what exactly determines those quantities (money demand and supply) to which the central bank reacts? And does it react so as to maximise its own surplus, to maximise welfare, to maximise the state's revenue, or a mix?

% }
% \mytext{
%         Where does the (long term) interest rate come from?
%         \begin{itemize}
%             \item Determined by markets?
%             \item Set by Lender of Last Resort?
%             \item To which end?
%             \item Determined by growth rate?
%         \end{itemize}
% }
% \end{frame}

% \begin{frame}{Nature of Interest Rates}
% \mytext{
%     {Determined by markets?}
%     \begin{itemize}
%         \item In Economics, everything is a market (in the long run)
%         \item Impossible to rule against markets
%         \item[] (Story of King Knuth, Diocletian)
%         \item Do we understand markets perfectly?
%         \item Can we model markets?
%         \item Simplification, model, map an territory        
%     \end{itemize}
% }
% \end{frame}
% \begin{frame}{Nature of Interest Rates}
% \mytext{
%     {Set by Lender of Last Resort?}
%     \begin{itemize}
%         \item Central Bank (CB, national central bank)
%         \item Leading central bank (Federal Open Market Committee, US)
%         \item IMF (Special Drawing Rights), BIS?
%     \end{itemize}

%     In any case:
%     \begin{itemize}
%         \item What is the policy of that LoLR $\rightarrow$ public CB
%         \item Price stability ($\pi=0$) (of which basket?) $\rightarrow$ independent CB
%         \item Own (real) income $\rightarrow$ private CB (US?)
%         \item Autonomous, or market participant?
%     \end{itemize}
% }
% \end{frame}

% \begin{frame}{Nature of Interest Rates}
% \mytext{
%     {Determined by Growth Rate?}
%     \begin{itemize}
%         \item What determines the growth rate $g$?
%         \item What is in the basket of $P$?
%         \item Technological progress?
%         \item Population growth?
%     \end{itemize}
% }    
% \end{frame}


\begin{frame}{Inflation}
    \mycomment{
    Inflation is a rise of the price level. This means that the monetary base $M$ increases faster than output $Y$. The main result is that wages, which are typically negotiated in larger intervals than prices, go down with inflation. This makes production cheaper.
    }
    \mytext{
    \begin{itemize}
        \item rising prices lowers real wages
        \item under full employment:
        \begin{itemize}
            \item unionisation etc.
            \item wage goes up
            \item demand goes up
            \item prices go up
            \item feedback loop
            \item CPI consumer price index
        \end{itemize}
        \item no full employment:
        \begin{itemize}
            \item more production
            \item more employment
            \item higher income
        \end{itemize}
    \end{itemize}    
    }
\end{frame}



\begin{frame}{Phillips Curve}
    \showmytikz{../shared/tikz/tikz-phillipscurve}
    \mycomment{
    The Phillips curve conveys the idea that there is a relationship between inflation and unemployment.

    This was generally considered true until the age of stagflation, where inflation and unemployment happen together.

    Supply shocks may shift this curve to the right; and it can generally be rather unstable. The reality seems more complex.
    }
\end{frame}

\begin{frame}{Circular Model of the Economy}
    \showmytikz[1,2,3]{../shared/tikz/tikz-flowmodel-own}
    \mycomment{
        In the context of macroeconomics, the economy is circular: households save (on capital markets), their savings become investments (of business). Households buy and consume (goods and services), which are produced by businesses.

    }
\end{frame}

\begin{frame}{Circular Model of the Economy}
    \showmytikz[1,2,3,4]{../shared/tikz/tikz-flowmodel-own}
    \mycomment{
        
        On top, households pay taxes (to the government), which the government spends on markets (to businesses). 
        
        Households also import goods and services from the rest of the world, while business export to the rest of the world.

        Microeconomics looks at individual branches of this picture: how households (consumers) decide, how businesses decide, and how the state (should) decide, albeit \textit{without taking macroeconomic effects into account.}

        If we understand the economy as a whole (macroeconomics) as a complex system of interrelated equations, we here (in microeconomics) deal with partial equations, and partial derivatives, and partial optimisation.

        % Note that this is simplified. The government may also redistribute, and has other functions as well. Households decide how much time (and skill) to offer on labour markets. There is also an international capital market. Businesses may be organised in one way or the other, and households may be large or small.
    }
\end{frame}


% \begin{frame}{Money Supply and Money Demand}
% \mytext{
%     \begin{itemize}
%         \item Demand for money is demand for credit
%         \item interest rate result of credit market
%         \item or: demand for money for economic activity
%         \item interest rate result of money market
%         \item disputed...
%     \end{itemize}
% }
% \end{frame}

\begin{frame}{IS-LM Model: IS Curve}
    \showmytikz[1,2]{../shared/tikz/tikz-islm}
    \mycomment{
   $IS$ is the curve of all interest rate - output - combinations that result in a balance of the real markets, that is, of investments~$I$ and savings~$S$.

    At $P_1$, a relatively high interest rate makes credit expensive. At this rate, much will be saved, hence less will be spent on consumption. There will be relatively low output.

    At $P_2$, lower interest rate make consumption more interesting than saving. Thus, output $Y$ will be higher.

    % Empirically, the $IS$ curve is relatively steep, while the $LM$ curve is relatively flat. 
    }
\end{frame}

\begin{frame}{IS-LM Model: LM Curve}
    \showmytikz[3,4]{../shared/tikz/tikz-islm}
    \mycomment{
        $LM$ is the curve of all combinations of interest rate and output that result in a balance on the money market, namely of liquidity preference $L$ and money supply $M$.

        At $P_3$, a relatively low output will require little money, and low money demand will result in low interest rates. 

        At $P_4$, a higher output will require more money, and as more profits are realised, interest rates will rise.

        The exact shape of the $LM$ curve is less clear than that of the $IS$ curve. In reality, it might be very flat. Some say it is flat on the left (liquidity trap) and steep on the right.        
        
        % Future expectations and policy play a larger role on the financial markets than in the real markets. 
    }
\end{frame}

\begin{frame}{IS-LM Model: Equilibrium}
    \showmytikz[1,3,5]{../shared/tikz/tikz-islm}
    \mycomment{
        There is an equilibrium $A$ where $IS$ and $LM$ meet, that is the synchronous balance of real and financial markets. At this combination of interest rate $i$ and output $Y$ investments will equal savings (the point is on $IS$) and money supply will equal money demand (the point is on $LM)$.
    }
\end{frame}


\begin{frame}{IS-LM Model: Disequilibria}
    \showmytikz[1,2,3,4,5]{../shared/tikz/tikz-islm}
    \mycomment{
        $P_1$: high interest rates result in little spending and low output. At this interest rate there is an excess of money supply (the point is left of $LM$). Money suppliers will lower the interest rate to find clients, which will result in a higher output, so movement towards $A$.

        $P_2$: at a low interest rate and high output, there is excess demand for money for transactions. This will drive up interest rates. 

        $P_3$: low interest rate and low output means there is much cheap money for more consumption available, creating excess demand for goods, driving up output. 

        $P_4$: high interest rates clear the money market, but the output is too high at this point for a balanced real market; subjects would rather save than consume this much. Thus, production and interest rates will have to go down.
        
    }
\end{frame}



\begin{frame}{IS-LM Model: Fiscal Policy}
    \showmytikz[1,3,5,6]{../shared/tikz/tikz-islm}
    \mycomment{
        Government expenditure can shift the $IS$ curve to the right, resulting in higher interest rate $i_2>i_1$ and larger output $Y_2>Y_1$. 
        
        The effect is only partial; the rise in output is smaller than the additional government expenditure. The risen interest rate has resulted a partial crowding-out. 

        Tariffs have a similar effect, as they raise net exports ($NX=EX-IM$). 

        Any fiscal policy has also distributive effects not shown here.
    }
\end{frame}

\begin{frame}{IS-LM Model: Monetary Policy}
    \showmytikz[1,3,5,7]{../shared/tikz/tikz-islm}
    \mycomment{
        Monetary policy (by the central bank) can shift the $LM$ curve to the right. This can happen either through lowering of the interest rate, or by open market policy. 

        The new equilibrium shows a lower interest rate and higher output $Y_3>Y_1$. The effect is also only partial: the rise in output is smaller than the monetary stimulus, since the lower interest rate has reduced savings and investment.

        There is an obvious contradiction (or: competition?) between fiscal and monetary policy, which is the reason why central banks \textit{should be} independent.
    }
\end{frame}

\begin{frame}{IS-LM Model: Monetary Policy}
    \showmytikz[1,3,5,6,7]{../shared/tikz/tikz-islm}
    \mycomment{
        Fiscal and monetary policy can complement or offset each other. Expansion by policy will meet its natural boundary at unsustainable government debt levels or at runaway inflation discrediting the currency. 
        
        It is thus widely believed that any such policy should be used temporarily only, like to overcome an acute crisis, external shocks, slumps, or maybe to catch up on development. 
        
        It is a bit like burning furniture to heat the room, or selling the silverware: you can do it in a crisis, but it is not a long-term solution.
    }
\end{frame}


    
\begin{frame}{IS-LM Model: Monetary Policy}
    \showmytikz[1,3,5,6,7]{../shared/tikz/tikz-islm-smooth}
    \mycomment{
        The actual shape of IS and LM are debated. They `seem' to be steep at the extremes, but flat for low levels of $i$. This can, of course, have dramatic consequences whether fiscal or monetary policy works at all. We will not discuss this here in detail.
    }
\end{frame}


    
